\documentclass[11pt,a4paper]{article}

\usepackage[margin=2.5cm]{geometry}
\usepackage{hyperref}
\usepackage{booktabs}
\usepackage{array}
\usepackage{parskip}
\usepackage{titlesec}
\usepackage{enumitem}
\usepackage{xcolor}

\hypersetup{
  colorlinks=true,
  linkcolor=blue,
  urlcolor=blue,
  citecolor=blue
}

\titleformat{\section}{\large\bfseries}{\thesection}{1em}{}
\titleformat{\subsection}{\normalsize\bfseries}{\thesubsection}{1em}{}

\title{\textbf{Cupel: A Framework for Professional Competence Provenance}\\
\large Verifying Human Expertise in a World of Capable Machines\\
\normalsize Position Paper and Framework Specification — v0.4}

\author{
  Eugene Andrie Merwe-Chartier\\
  \href{https://cupel.foundation}{cupel.foundation}\\
  \href{mailto:contact@cupel.foundation}{contact@cupel.foundation}
}

\date{July 2026 (v0.4; previous versions: v0.1 11 March 2026, v0.2 18 March 2026, v0.3 16 May 2026)}

\begin{document}

\maketitle

\begin{abstract}
The deployment of autonomous and semi-autonomous artificial intelligence agents into professional
workflows has created a structural verification problem in the human layer those systems depend on.
Two trends are converging. First, AI systems can now produce the signals of professional competence
— passing licensing examinations, generating credible portfolios, drafting credentialed work — without
possessing the judgment, accountability, or contextual experience that credentials were designed to
indicate. Second, the apprenticeship pipeline through which qualified human supervisors are produced
is contracting: entry-level technology hiring at the top fifteen firms fell 25\% between 2023 and
2024~\cite{signalfire2024}, and UK graduate technology roles fell 46\% in 2024~\cite{rezi2026}. These
trends compound: organisations are deploying systems faster than they can identify the humans
capable of supervising them, while simultaneously narrowing the pipeline that produces such humans.
This paper identifies the resulting gap, characterises the property of competence that matters in
agent-mediated work — what we term competence \emph{below the abstraction layer} — and proposes
Cupel, an open framework specification for professional competence provenance. Cupel defines a
taxonomy of five trust signal types, a JSON-LD schema for representing and linking those signals,
and an interoperability architecture that sits above existing standards including W3C Verifiable
Credentials, C2PA, the Credential Transparency Description Language (CTDL), and Open Badges 3.0.
This revision adds an explicit treatment of the human stakes behind the pipeline data: early 2026
labour-market evidence shows the entry ramp into knowledge work narrowing before any broad
unemployment signal appears, which changes what an adequate policy and infrastructure response
looks like. The framework is released under Apache 2.0, with trademark protection ensuring
conformance integrity.
\end{abstract}

\tableofcontents
\newpage

%% ─────────────────────────────────────────────────────────────────────────────
\section{Introduction}

When an autonomous agent files a tax return, drafts a contract, recommends a clinical course of
action, or commits production code, the integrity of that output depends on a qualified human who
can detect its errors, intervene meaningfully, and accept professional responsibility. The
credential system that organised professional labour markets for most of the twentieth century was
designed in part to identify those humans: a chartered accountant, licensed physician, or admitted
barrister was not merely a knowledge marker but an accountability node embedded in a system of
professional bodies, regulatory oversight, and liability mechanisms.

This identification system is failing on two fronts simultaneously, and the failure is
accelerating.

\textbf{The signal is no longer informative.} The effectiveness of credentials as competence proxies
depended on a single equilibrium condition: \emph{gaming the proxy must be roughly as hard as
becoming genuinely competent}. This condition has been violated. In September 2025, large language
models passed CFA Level III mock examinations including written essays, with leading models scoring
above 79\%~\cite{nyustern2025}. The same pattern has emerged across the bar exam, medical licensing
examinations, and engineering certification assessments. AI systems can acquire the signal without
acquiring the substance. This is a textbook instantiation of Goodhart's Law: when a measure becomes
a target, it ceases to be a good measure~\cite{goodhart1975}. The market response — issuing more
credentials, from 334,114 unique credentials in the United States in 2018 to 1.85 million by
2025~\cite{credentialengine2025} — has increased signal volume while continuing to degrade signal
quality. Gartner projects that one in four candidate profiles worldwide will be entirely fabricated
by 2028, and 41\% of enterprises report having already onboarded a fraudulent
candidate~\cite{gartner2025}.

\textbf{The pipeline that produces qualified supervisors is contracting.} Entry-level technology
hiring at the top fifteen firms fell 25\% between 2023 and 2024~\cite{signalfire2024}; UK graduate
technology roles fell 46\% in 2024~\cite{rezi2026}. Junior roles are not interchangeable with senior
ones — they are the structured apprenticeship layer through which senior practitioners are formed.
Organisations compressing junior hiring are defunding their own future supervisory capacity. The
problem is time-delayed: by the time the shortage of qualified supervisors becomes acute, the cohort
that would have filled those roles will not exist.

\textbf{Deployment is outpacing oversight.} Multi-agent system adoption rose 327\% in four
months~\cite{databricks2026}. Only one in five companies has mature governance for autonomous
agents~\cite{deloitte2026}. 78\% of AI users bring personal tools to work, often outside any
oversight framework~\cite{aigov2025}. The European Union's AI Act activates high-risk compliance
obligations through 2026, with Article 14 requiring meaningful human oversight and Article 50
imposing content provenance obligations~\cite{euaiact2024}. These obligations presuppose a
verifiable population of competent supervisors. No mechanism currently exists to identify that
population at scale.

The problem is structural, not cyclical. The signal layer of professional labour markets and the
supervisory layer for AI deployment have become the same problem: organisations need to know which
humans can genuinely catch the errors that machines now produce at scale. This paper describes the
infrastructure that closes that gap: the Cupel Professional Competence Provenance Framework.

%% ─────────────────────────────────────────────────────────────────────────────
\section{Below the Abstraction Layer}

A central claim of this paper is that the property of competence which now matters in
agent-mediated work is not well captured by existing credentialing frameworks. We call this
property competence \emph{below the abstraction layer}.

Most professional credentials measure performance \emph{above} the AI abstraction layer — on
tasks that AI systems can now also perform. A candidate who can pass the CFA Level III
demonstrates an ability that a large language model also demonstrates. The credential, on its own,
no longer discriminates between the two. Performance at this level has become diagnostic of
nothing more specific than the ability to produce credentialed output.

Competence below the abstraction layer is qualitatively different. It is the ability to:

\begin{itemize}[noitemsep]
  \item recognise when an AI output is wrong, in cases where the error is non-obvious to a
        non-expert evaluator;
  \item identify the specific dimension on which the output is wrong (factually, structurally,
        contextually, ethically);
  \item determine what corrective action is appropriate and, where applicable, execute it;
  \item accept professional responsibility for the corrected output under the relevant
        accountability framework.
\end{itemize}

Three illustrative cases clarify the distinction:

A clinician using an AI system to read a diagnostic scan is not abdicating responsibility to the
system. The clinician retains diagnostic authority. The signal that matters is whether the
clinician can recognise a subtle finding that the model missed or misclassified — not whether the
clinician could replicate the model's general performance on a licensing examination.

A code reviewer approving an AI-generated patch must be able to identify a security vulnerability
that the model did not flag. The reviewer's competence is judged by what they catch, not by what
they would produce unassisted.

A compliance officer signing off on an AI-drafted disclosure must understand the regulatory rule
the AI applied, the conditions under which that rule has exceptions, and whether the present case
falls within those exceptions. The signature is meaningful only if the officer has the capacity
to dissent.

In each case, the human is exercising what regulators term ``human oversight'' or ``human in the
loop''~\cite{euaiact2024}. The legal and ethical force of that oversight depends entirely on
whether the human can in fact catch the loop's mistakes — a property that traditional credentials
were not designed to verify, and that AI-era credential gaming makes harder to verify by surface
inspection.

Cupel's design centres on producing legible signals of competence below the abstraction layer.
The five trust signal types defined in Section~5 are each ways of evidencing this property; the
multi-signal model is the framework's principal answer to the question of how such competence
can be verified at scale without recreating the gameable conditions of single-signal credentialing.

%% ─────────────────────────────────────────────────────────────────────────────
\section{The Human Stakes}

The pipeline data cited in Section~1 is not merely a supply-side constraint on future supervisory
capacity; it describes a present and specific labour-market anxiety, and the distinction matters
for what an adequate response looks like.

Anthropic's Economic Index reports that, as of early 2026, hiring of 22–25 year-olds into the
occupations most exposed to AI has slowed by approximately 14\%, while no clear broad unemployment
signal has yet appeared in those same occupations~\cite{aei2026}. Roughly one in five workers in
AI-exposed jobs report concern about economic displacement, with software engineers among the most
concerned of any profession surveyed~\cite{aei2026}. Read together, these findings indicate that
the entry ramp into knowledge work is narrowing \emph{before} any headline unemployment figure
moves, because employers can defer hiring decisions long before they take the more visible step of
reducing headcount.

This distinction is consequential. If the operative crisis were simply contracting employment, the
appropriate policy response would be conventional: income support and retraining. But if the
operative crisis is that the apprenticeship structure converting novices into accountable experts is
quietly failing — junior roles being the mechanism by which practitioners historically acquired the
below-the-abstraction-layer competence characterised in Section~2 — then income support and
retraining do not address the underlying failure. What is additionally required is verification
infrastructure that can recognise competence in practitioners who did not pass through that
structure in the traditional way, and that can make legible who, among a shrinking cohort, is
actually capable of supervising agent-mediated work.

This also reframes who has standing in the conversation this paper is part of. The verification
problem is not confined to standards bodies, credentialing vendors, and enterprise HR functions. It
is a live concern for anyone whose career ladder has had rungs removed, and for anyone deciding how
much unsupervised authority to extend to a system that no one has yet verified a human can
adequately oversee. A parallel and growing concern, documented separately, is that AI has made
confident commentary about these dynamics cheap to produce and correspondingly easy to overtrust,
widening the gap between people describing the future of knowledge work and people building a
working response to it~\cite{hbr2026}. This paper, and the open framework it proposes, is offered as
evidence-based participation in that response rather than commentary on it.

%% ─────────────────────────────────────────────────────────────────────────────
\section{What Credentials Actually Did}

Understanding what must replace credentials requires clarity on what they actually provided. Four
distinct functions are now unbundling:

\textbf{Information asymmetry resolution.} Credentials were trust shortcuts for principals who
cannot evaluate agent competence directly. If AI can produce competent-seeming professional output,
the principal now faces \emph{worse} information asymmetry: they cannot determine whether the
professional is applying judgment on top of AI output or simply relaying it.

\textbf{Accountability infrastructure.} Professional credentials were embedded in systems of licence
revocation, malpractice liability, and disciplinary action. AI-generated output carries no equivalent
accountability infrastructure. The credential was the entry condition for accountability, not merely
a knowledge signal.

\textbf{Community of practice gatekeeping.} Much professional expertise is tacit knowledge
distributed within practice communities. The credential was the admission ticket. This gatekeeping
function persists even as the knowledge-signalling function collapses, creating an increasingly
hollow dynamic.

\textbf{Scalable evaluation.} Credential verification scales: a recruiter can screen 10,000
applicants by checking degrees. Any replacement system must address scalability. Most proposed
alternatives — portfolio review, work samples, structured reference checks — are high-bandwidth
but low-scale. This asymmetry is the central engineering challenge.

%% ─────────────────────────────────────────────────────────────────────────────
\section{Emerging Trust Mechanisms}

Five new trust mechanisms are forming at different stages of maturity. The Cupel framework is
designed to provide infrastructure for all five.

\subsection{Demonstrated Outcome Trails}
Verifiable records of actions taken and outcomes achieved: portfolios, open-source contributions,
competition results, documented client outcomes, and public track records. This is the most mature
emerging mechanism. Its limitation is domain specificity — outcome trails that are legible in
software engineering (GitHub) or data science (Kaggle) have no equivalent in consulting,
legal practice, or financial advisory.

\subsection{Behavioural Consistency Signals}
Consistency across time and contexts as a primary trust signal. A professional with a decade of
context-appropriate judgment documented across engagements signals something a certification cannot.
Consistency is hard to fake because it requires sustained performance. It is also hard to
\emph{capture} without infrastructure purpose-built for that function.

\subsection{Network-Verified Reputation}
Professional networks becoming the primary trust mechanism rather than a supplement. 63\% of
consumers identify verified reviews as the most important trust signal~\cite{brightlocal2025}. In
professional contexts this manifests as referral networks, public endorsements, and the
``who vouches for you'' graph. The limitation is susceptibility to network homophily and
collusion.

\subsection{Process Transparency and Human-AI Collaboration Records}
The new credential is not ``I know finance'' but ``here is how I use AI in my process, here is
where I override it, and here is my track record of those overrides being correct.'' Trust shifts
from \emph{what you know} to \emph{how you think} alongside AI systems. This is nascent but
structurally important as AI capability continues to advance.

A critical property of this signal type is that it captures \emph{attribution clarity} — not a
ratio of human versus machine effort, but a record of at which decision points human judgment was
exercised, what alternatives were available, and what the consequences were. A practitioner who
delegates execution to AI while making all structural choices is exercising full professional
accountability; the relevant signal is the quality and track record of those structural choices,
not the volume of manual work.

\subsection{Provenance and Attribution Infrastructure}
Content credentials, human-in-the-loop documentation, and chain-of-custody metadata for
professional work product. The term \emph{content credentials} follows C2PA's established
nomenclature for cryptographically-bound provenance metadata attached to a digital
artefact~\cite{c2pa2024}. The analogy is organic food labelling: a mechanism for distinguishing
human-exercised judgment from AI-generated output at the level of individual work products.
This layer is being accelerated by regulatory requirements including EU AI Act Article 50
content provenance obligations and the California AI Transparency Act (disclosure requirements
effective January 2026).

Critically, provenance-tracked signals must capture an \emph{AI audit trail} — a structured
log of process metadata — in addition to final outputs. The term ``AI audit trail'' is now
established in both regulatory compliance contexts and enterprise AI governance
frameworks~\cite{audit2025}. An AI system can produce a polished final deliverable that is
indistinguishable from expert human output. What AI cannot replicate authentically is the
iterative trail of a human-led process: the drafts, corrections, mid-course reversals, and
contextually-situated judgments that characterise genuine professional engagement with a problem.
A provenance-tracked signal for a piece of professional work should therefore include, at minimum:
the content credential (which AI systems, if any, were used and in what capacity, expressed in
a C2PA-compatible format), the AI audit trail (evidence of human decision points across the
development of the work), and the accountability anchor (who is responsible for the outcome and
under what professional framework).

\subsection{The Structural Tension}

Old credentials were \emph{scalable but low-bandwidth}. New trust mechanisms are
\emph{high-bandwidth but low-scale}. This asymmetry creates market bifurcation: high-value
professional work will operate on rich, verified trust signals; routine work will be commoditised
into AI pipelines with binary quality checks. The middle layer of professional practice faces
the most disruption.

%% ─────────────────────────────────────────────────────────────────────────────
\section{Market Landscape and the Gap}

The current market for professional trust infrastructure can be characterised as five layers,
each with distinct players and commercial dynamics.

\begin{table}[h]
\centering
\small
\begin{tabular}{@{}p{0.22\textwidth}p{0.35\textwidth}p{0.35\textwidth}@{}}
\toprule
\textbf{Layer} & \textbf{Representative Players} & \textbf{Maturity} \\
\midrule
Identity Verification \& Fraud Detection
  & Yardstik, GetReal Security, Pindrop, Sherlock AI
  & Most funded; pain-driven buying \\
\addlinespace
Skills Assessment \& Validation
  & TestGorilla, HireVue, Codility, Pymetrics
  & Growing fast; 41\% of firms now use assessments \\
\addlinespace
Digital Credential Infrastructure
  & Credential Engine (CTDL), MIT DCC, 1EdTech (CLR)
  & Standards-building; technically strong, commercially early \\
\addlinespace
Content Provenance \& Authenticity
  & C2PA coalition (200+ members inc.\ Adobe, Microsoft, Google)
  & Regulatory-driven; ISO fast-track underway \\
\addlinespace
Professional Reputation \& Outcome Tracking
  & LinkedIn (self-reported), Crosschq, Toptal, Kaggle
  & Least developed; biggest structural gap \\
\bottomrule
\end{tabular}
\caption{The five-layer professional trust infrastructure landscape.}
\end{table}

Most investment is concentrated in Layers 1 and 2: detecting fraud and testing skills. These are
defensive plays. The \emph{constructive} infrastructure — systems that positively identify
competence and track demonstrated judgment over time — is severely underdeveloped.

The critical absence is a \textbf{cross-cutting evaluation and interoperability layer}: a common
language for describing, linking, and evaluating professional competence signals across the entire
fragmented ecosystem. With 1.85 million credentials in circulation and widespread inability to
evaluate them, this aggregation layer represents the most significant unmet structural need in
professional labour markets.

The logical convergence point is a unified professional competence provenance system that
connects all five layers into a coherent, interoperable whole. Whoever defines the architecture
for this convergence is building the replacement for the degree.

%% ─────────────────────────────────────────────────────────────────────────────
\section{The Cupel Framework}

\subsection{Design Rationale: Open Standard Over Product}

The Coalition for Content Provenance and Authenticity (C2PA) provides the most instructive
precedent. C2PA succeeded not by launching a product but by publishing a specification and building
a coalition. It is now fast-tracked as an ISO standard and has over 200 member organisations.
The professional competence provenance problem is structurally similar: it requires cross-industry
adoption, trust from multiple competing stakeholders, and interoperability between platforms with
conflicting commercial interests. An open standard maximises adoption and positions the framework
as neutral infrastructure rather than as a competitor to existing players.

\subsection{Framework Components}

Cupel specifies five interrelated components:

\textbf{Trust Signal Taxonomy.} A structured classification of professional competence signals
into five types:
\begin{itemize}[noitemsep]
  \item \emph{Credential-based}: issued qualifications, licences, certifications
  \item \emph{Assessment-based}: structured evaluations of demonstrated skill
  \item \emph{Outcome-based}: verifiable records of actions and consequences
  \item \emph{Peer-verified}: endorsements and attestations from identified third parties
  \item \emph{Provenance-tracked}: content credentials and AI audit trails for professional
        work product, capturing both chain-of-custody metadata and the process record
        (AI tool disclosure, human decision points, iteration history)
\end{itemize}

The five-type model is designed for signal \emph{diversity} — what practitioners in
skills-based hiring increasingly term a \emph{trust portfolio}~\cite{sbhiring2025}: no single
signal type is sufficient, and the evidential weight of a competence claim increases with
corroboration across multiple independent signal types. This design makes the framework resistant
to single-vector gaming: inflating peer-verified signals, for example, does not compensate for
absence of outcome-based or provenance-tracked evidence.

\textbf{Core Schema.} A JSON-LD schema for expressing trust signals and their relationships,
aligned with the W3C Verifiable Credentials Data Model v2.0~\cite{w3cvc2022}. The schema
captures signal type, issuer identity, subject identity, temporal scope, domain specificity,
and verification method. For provenance-tracked signals specifically, the schema includes a
\texttt{contentCredential} field (C2PA-compatible declaration of AI tools used and their scope
of involvement) and an \texttt{auditTrail} field (structured log of human decision points and
process iterations). These fields are optional for other signal types but required for
provenance-tracked signals.

\textbf{Selective Disclosure Principle.} The framework defines a vocabulary for describing
competence signals; it does not mandate a surveillance architecture. Conformant implementations
must support selective disclosure: a practitioner controls which signals are shared, with which
parties, and for which purposes. This principle is consistent with W3C Verifiable Credentials'
selective disclosure capabilities and with the EU General Data Protection Regulation's data
minimisation requirements. Provenance tracking creates a potential privacy tension — detailed
process logs could expose proprietary methods or working patterns — and the selective disclosure
principle is the framework's structural answer to that tension.

\textbf{Signal Quality Criteria.} Evaluation dimensions for assessing the evidential weight of
individual signals: source credibility, recency, verifiability, specificity, and consistency
with concurrent signals.

\textbf{Interoperability Architecture.} Defined mappings between Cupel signal types and
existing standards, enabling bidirectional translation without requiring issuing platforms to
adopt the full framework.

\textbf{Governance Model.} A trust list maintenance mechanism, conformance programme, and
contributor framework modelled on established open standards governance (W3C, IETF, C2PA).

The governance model must specifically address what may be termed the \emph{oracle problem of
attribution}: if Cupel relies on issuing systems (certification bodies, HR platforms, assessment
providers) to supply signal metadata, the framework's trust guarantees are bounded by the
integrity of those systems. A dishonest or negligent issuer can supply structurally conformant
but substantively false signals.

The framework addresses this through an \emph{issuer trust registry}: a reputation mechanism
in which issuers accumulate credibility based on the observable real-world correlation between
their signals and verified outcomes over time. The term ``trust registry'' is established in
W3C Verifiable Credentials contexts and in MIT's Digital Credentials Consortium Governance
Framework for Issuer Identity Registries (June 2025)~\cite{mitdcc2025}. This is a market
mechanism, not a technical mandate. It follows the pattern of email sender reputation systems
(SPF/DKIM scoring) and OpenID Federation's trust chains: no central authority assigns
credibility; credibility emerges from the longitudinal track record of each issuer's signals.
Evaluators weight signals from high-integrity issuers more heavily, and the market incentivises
issuers to maintain signal quality. The registry itself requires ongoing data collection and
curation infrastructure; the specification defines the interface but does not mandate a single
operator.

\subsection{Relationship to Existing Standards}

Cupel does not replace existing standards. It sits above them as an evaluation and
interoperability layer:

\begin{table}[h]
\centering
\small
\begin{tabular}{@{}p{0.28\textwidth}p{0.65\textwidth}@{}}
\toprule
\textbf{Standard} & \textbf{Role in Cupel Architecture} \\
\midrule
W3C Verifiable Credentials & Cryptographic verification layer for all signal types \\
C2PA & Content provenance architecture; template for work product provenance signals \\
Credential Engine / CTDL & Structured description of credential-based signals \\
Open Badges 3.0 / CLR & Micro-credential and learner record formats for assessment-based signals \\
ISO/IEC 27001, 29101 & Security and privacy frameworks for implementation conformance \\
\bottomrule
\end{tabular}
\caption{Cupel's relationship to existing standards.}
\end{table}

The novel contribution is the cross-cutting layer: a unified representation that allows a
composite professional competence provenance profile to be constructed from signals originating
in any of the above ecosystems.

%% ─────────────────────────────────────────────────────────────────────────────
\section{Licensing and Conformance}

The framework specification and core schema are released under Apache 2.0. Any implementation
may use the technical content without restriction. The Cupel name and logo are trademarked
(UK00004352899, filed 11 March 2026), ensuring that conformance claims carry meaning: only
implementations that meet published conformance criteria may identify themselves as Cupel
implementations. This follows the Linux and OpenID precedent, where an open technical artefact
is paired with a protected name to prevent the dilution of conformance through unverified
implementations.

The selection of Apache 2.0 in preference to a copyleft licence reflects an adoption-first
strategy. The framework's value depends on broad cross-industry uptake; a permissive licence
removes friction for commercial implementers and standards bodies that would be unable to adopt
a strong copyleft licence as a matter of internal policy. The specification is intended to be
usable without restriction by both open and proprietary implementations.

Specific operational infrastructure that supports the framework — for example, an issuer trust
registry that aggregates real-world signal-outcome correlations — may be operated by independent
parties under their own governance arrangements. The specification defines what such
infrastructure must do to be conformant; it does not nominate a single operator and does not
preclude any particular funding model.

%% ─────────────────────────────────────────────────────────────────────────────
\section{Conclusion}

The credential system that organised professional labour markets for most of the twentieth century
is structurally compromised, and the supervisory layer that AI deployment depends upon rests on
the same compromised foundation. The cause is not a failure of operators but the convergence of
three trends: AI systems that can produce the surface signals of competence without the
substance, a contracting apprenticeship pipeline that is narrowing the cohort of practitioners
from which future supervisors will be drawn, and an accelerating deployment of agents into roles
that presuppose competent human oversight.

The market response to date has concentrated on defensive infrastructure: detecting fraud,
testing skills in controlled environments, verifying that AI outputs meet quality thresholds.
The constructive work of building new positive trust infrastructure — systems that identify and
track the humans capable of supervising agent work over time — is largely undone.

Cupel proposes the architecture for that infrastructure. It is released as an open standard
because the problem requires coalition, not competition. The framework's contribution is
deliberately bounded: it specifies a common vocabulary, a data format, and an interoperability
architecture. It does not assert what makes someone competent; it provides the substrate on
which competence claims from many sources can be expressed, linked, and evaluated together. The
specification, reference implementation in development, and governance framework are available
at \href{https://github.com/vdmeu/CUPEL}{github.com/vdmeu/CUPEL} and
\href{https://cupel.foundation}{cupel.foundation}.

The framework's adoption strategy targets the alternative credential and micro-assessment
ecosystem as the initial cohort: coding bootcamps, AI-tutor platforms, portfolio-based
assessment providers, and open-source contribution tracking tools. These organisations have
the most to gain from interoperability — their signals are currently unreadable to enterprise
hiring systems — and no incumbent lock-in to protect. Broad adoption in the long tail of the
credential ecosystem creates the interoperability pressure that eventually compels larger
platforms to participate, mirroring the bottom-up adoption dynamics of C2PA and Open Badges.
Regulatory tailwinds — particularly the EU AI Act's transparency and oversight requirements —
provide additional leverage with larger incumbents for whom compliance creates a structural
incentive to adopt standardised provenance formats.

Several open questions remain, and the framework is offered for community engagement rather
than as a finished proposal. Among them: the technical specification of a live provenance signal
data model; the operational design of selective disclosure under the General Data Protection
Regulation and equivalent regimes; the resistance properties of the multi-signal model against
sophisticated gaming strategies; and the appropriate governance arrangements for an issuer
trust registry. Community input on the taxonomy, schema, and interoperability mappings is
welcomed through the repository's issue tracker and discussion board.

%% ─────────────────────────────────────────────────────────────────────────────
\section*{About the Author}

Eugene Andrie Merwe-Chartier writes and maintains Cupel as an open framework and an ongoing,
evidence-based public argument, not as a commercial product. Correspondence, critique, and
contributions are welcome at \href{mailto:contact@cupel.foundation}{contact@cupel.foundation} and
\href{https://github.com/vdmeu/CUPEL}{github.com/vdmeu/CUPEL}, where periodic commentary tracking
the framework's three constituent threads — provenance, credentialing, and the human impact of
automated knowledge work — is also published.

\section*{Acknowledgements}

The analytical foundation for this framework was developed in February 2026 and is documented
in full at \href{https://github.com/vdmeu/CUPEL/blob/main/docs/00-genesis/conversation-context.md}{docs/00-genesis/conversation-context.md},
which serves as the timestamped record of the intellectual origin of this work.

%% ─────────────────────────────────────────────────────────────────────────────
\begin{thebibliography}{9}

\bibitem{nyustern2025}
NYU Stern School of Business and Goodfin (2025).
\textit{AI Performance on the CFA Examinations}.
Working Paper, NYU Stern School of Business.

\bibitem{signalfire2024}
SignalFire (2024).
\textit{State of Tech Talent Report: Entry-Level Hiring Trends at the Top 15 Tech Firms}.
SignalFire Research.

\bibitem{rezi2026}
Rezi.ai (2026).
\textit{UK Tech Graduate Roles Report: 2024 Hiring Outcomes}.
Rezi.ai, January 2026.

\bibitem{deloitte2026}
Deloitte (2026).
\textit{State of Generative AI in the Enterprise: Governance and Risk}.
Deloitte Insights.

\bibitem{databricks2026}
Databricks (2026).
\textit{Multi-Agent System Adoption Report}.
Databricks Research.

\bibitem{aigov2025}
AI Governance Benchmark (2025).
\textit{Shadow AI in the Workplace: Adoption and Oversight Gaps}.
AI Governance Benchmark Consortium.

\bibitem{goodhart1975}
Goodhart, C.A.E. (1975).
\textit{Problems of Monetary Management: The U.K. Experience}.
Papers in Monetary Economics, Reserve Bank of Australia.

\bibitem{credentialengine2025}
Credential Engine (2025).
\textit{Counting U.S. Postsecondary and Secondary Credentials}.
Washington, DC: Credential Engine.
\url{https://credentialengine.org/counting-credentials/}

\bibitem{gartner2025}
Gartner (2025).
\textit{Predicts 2025: Artificial Intelligence in Human Capital Management}.
Gartner Research.

\bibitem{shrm2025}
Society for Human Resource Management (2025).
\textit{Skills-Based Hiring and the Changing Value of Credentials}.
SHRM Research.

\bibitem{brightlocal2025}
BrightLocal (2025).
\textit{Local Consumer Review Survey 2025}.
\url{https://www.brightlocal.com/research/local-consumer-review-survey/}

\bibitem{c2pa2024}
Coalition for Content Provenance and Authenticity (2024).
\textit{C2PA Technical Specification v2.1}.
\url{https://c2pa.org/specifications/}

\bibitem{w3cvc2022}
Sporny, M., Longley, D., Chadwick, D. (2022).
\textit{Verifiable Credentials Data Model v1.1}.
W3C Recommendation. \url{https://www.w3.org/TR/vc-data-model/}

\bibitem{openbadges2023}
1EdTech Consortium (2023).
\textit{Open Badges Specification v3.0}.
\url{https://www.imsglobal.org/spec/ob/v3p0/}

\bibitem{audit2025}
Loe, I. (2025).
\textit{Your AI Agent Needs an Audit Trail, Not Just a Guardrail}.
Medium / AI Governance. \url{https://medium.com/@ianloe/your-ai-agent-needs-an-audit-trail-not-just-a-guardrail-6a41de67ae75}

\bibitem{mitdcc2025}
Digital Credentials Consortium, MIT (2025).
\textit{Governance Framework for Issuer Identity Registries}.
\url{https://digitalcredentials.mit.edu/docs/Governance-Framework-for-Issuer-Identity-Registries.pdf}

\bibitem{sbhiring2025}
National Association of Colleges and Employers (2026).
\textit{Skills-Based Hiring Grows but College Students Don't Fully Understand It}.
NACE. \url{https://www.naceweb.org/about-us/press/2026/skills-based-hiring-grows-but-college-students-dont-fully-understand-it}

\bibitem{euaiact2024}
European Parliament and Council (2024).
\textit{Regulation (EU) 2024/1689 on Artificial Intelligence (AI Act)}.
\url{https://artificialintelligenceact.eu/article/14/}

\bibitem{aei2026}
Anthropic (2026).
\textit{The Anthropic Economic Index Report}.
Anthropic Research, March 2026. \url{https://www.anthropic.com/research/economic-index-march-2026-report}

\bibitem{hbr2026}
Harvard Business Review (2026).
\textit{Has AI Ended Thought Leadership?}
HBR, March 2026. \url{https://hbr.org/2026/03/has-ai-ended-thought-leadership}

\end{thebibliography}

\end{document}
